% !TITLE 临江仙·梦后楼台高锁
% !AUTHOR 晏几道
% !DYNASTY 两宋
% !GENRE 词
% !ORDER 2

\begin{poem}
梦后楼台高锁，酒醒帘幕低垂。去年春恨却来时\textsuperscript{1}。
落花人独立，微雨燕双飞\textsuperscript{2}。

记得小蘋\textsuperscript{3}初见，两重心字罗衣\textsuperscript{4}。琵琶弦上说相思\textsuperscript{5}。
当时明月在，曾照彩云归\textsuperscript{6}。
\end{poem}

\begin{annotations}
\notetext{1}{却来时：又来到的时候。去年春天的离恨，在梦回酒醒时又一次涌上心头。}
\notetext{2}{落花人独立，微雨燕双飞：这两句是千古名句，化用自五代翁弘。花在落，人孤零零地站着；雨在飘，燕子却成双成对地飞——人与燕的“单”与“双”构成强烈反差。}
\notetext{3}{小蘋：晏几道友人家的歌女，是词人的心上人。这首词就是回忆与她分别后的思念。}
\notetext{4}{心字罗衣：领口绣有心字花纹的丝衣。一说“心字”指衣领形状像心形，一说指衣上的盘扣打成了心字结。两重心字——两人心意相通，心心相印。}
\notetext{5}{琵琶弦上说相思：在琵琶弦上弹出相思之情。小蘋善弹琵琶，通过音乐传达心意。}
\notetext{6}{彩云归：彩云飘回。彩云既指天上的云彩，也暗喻小蘋（像彩云一样美丽）。明月还是当时的明月，曾经照着彩云般的小蘋离去——如今月依旧，人已不见。}
\end{annotations}

\begin{appreciation}
上一首《浣溪沙》的作者晏殊，是这首词的作者晏几道的父亲。父子俩并称"大小晏"，词史上都有一号，可两个人的词是两个味道：晏殊是宰相，词里是体面的叹息；晏几道是落魄公子，词里是刻骨的怀念。这首《临江仙》是他最有名的作品之一，写的是一个忘不掉的人。

"梦后楼台高锁，酒醒帘幕低垂。"梦醒了，楼台的门锁着；酒醒了，帘幕垂着。晏几道这一辈子，好像总在梦后和酒醒之间——因为清醒的时候，回忆太重。"去年春恨却来时"，去年的离愁，今年不请自来。

"落花人独立，微雨燕双飞。"千古名句。花在落，人一个人站着；雨在飘，燕子成双成对地飞。一个"单"，一个"双"，孤独就藏不住了。这句其实不是他原创，是从五代诗人翁宏那里借来的——翁宏的原诗早没人记得，晏几道一用，就流传了一千年。一句诗落在谁的手里，也是缘分。

下阕写初见。"记得小蘋初见，两重心字罗衣。"小蘋是他家里的歌女。他记得她衣服上绣着两个心字，好像两人心意相通；她不会说相思，只在琵琶弦上弹出来——"琵琶弦上说相思"，用音乐传情，比说话还动人。"当时明月在，曾照彩云归。"当时的月亮还在天上，照着彩云一样的小蘋回去。月亮年年有，彩云般的人，再没见过。

书里《行行重行行》说"思君令人老"，晏几道大概是最懂这句话的词人之一。黄庭坚说他身上有"四痴"，其中一痴是：被人辜负一百次，也不记恨，还是照样相信人。这样的人写词，词里全是真话。

哥哥选这首词，是盼你做一个记得住美好的人。晏几道记了一辈子的是心字罗衣和琵琶声，你以后会记住的，是毕业照、同桌、大学宿舍——这些都值得记。至于伤心的事，忘了就好。晏几道太痴，你比他聪明。
\end{appreciation}
